% Part: computability
% Chapter: recursive-functions
% Section: primes

\documentclass[../../../include/open-logic-section]{subfiles}

\begin{document}

\olfileid{cmp}{rec}{pri}
\olsection{মৌলিক সংখ্যা}

সীমাবদ্ধ পরিমাণায়ন ও সীমাবদ্ধ ন্যূনীকরণ আমাদের এমন অনেক উপায় দেয়,
যেগুলির সাহায্যে স্বাভাবিক অপেক্ষক ও সম্বন্ধকে আদিম পুনরাবৃত্ত দেখানো যায়।
যেমন, “$x$ দ্বারা $y$ নিঃশেষে বিভাজ্য” সম্বন্ধটি বিবেচনা করুন; একে
$x \mid y$ লেখা হয়। $x \mid y$ সম্বন্ধটি সত্য হয় যদি $y$-কে~$x$
দিয়ে ভাগ করলে কোনো ভাগশেষ না থাকে, অর্থাৎ $y$ যদি~$x$-এর একটি
পূর্ণসংখ্যা গুণিতক হয়। (সম্বন্ধটি সত্য না হলে, অর্থাৎ $y$-কে $x$ দিয়ে
ভাগ করার ভাগশেষ $> 0$ হলে, আমরা $x \nmid y$ লিখি।) অন্যভাবে বললে,
$x \mid y$ যদি এবং কেবল যদি কোনো~$z$-এর জন্য $x \cdot z = y$। এমন
$z$ থাকলে স্পষ্টতই তা অবশ্যই $\leq y$। অতএব $x \mid y$ যদি এবং কেবল যদি কোনো $z \le y$-এর জন্য
$x \cdot z = y$। $=$ ও গুণ থেকে সীমাবদ্ধ অস্তিত্বমূলক পরিমাণায়নের
সাহায্যে $x \mid y$ সম্বন্ধটিকে এভাবে সংজ্ঞায়িত করা যায়:
\[
x \mid y \defiff \bexists{z \leq y}{(x \cdot z) = y}.
\]
সুতরাং আমরা দেখালাম যে $x \mid y$ আদিম পুনরাবৃত্ত।
% BN-SRC-141 TARGET-CORRECTION-BEGIN
\par\smallskip\noindent\textnormal{\small\textbf{উৎস-সংশোধন (BN-SRC-141):} হিমায়িত ইংরেজি পাঠে $x \nmid y$ ব্যাখ্যা করতে ভুল ক্রমে $x$-কে $y$ দিয়ে ভাগ করার কথা বলা হয়েছে। সংজ্ঞা $x \mid y$ অনুসারে বাংলা পাঠে $y$-কে $x$ দিয়ে ভাগ করার ভাগশেষ লেখা হয়েছে।} % BN-SRC-141 TARGET-CORRECTION-END

কোনো স্বাভাবিক সংখ্যা~$x$-কে
\emph{মৌলিক} বলা হয় যদি সেটি $0$ বা $1$ না হয় এবং কেবল $1$ ও নিজে
দ্বারা নিঃশেষে বিভাজ্য হয়। অন্যভাবে বললে, মৌলিক সংখ্যাগুলির ক্ষেত্রে
$y \mid x$ হলেই হয় $y = 1$, নয়তো~$y=x$। $x$ মৌলিক কি না পরীক্ষা
করতে আমাদের শুধু সব $y \le x$-এর জন্য $y \mid x$ কি না দেখতে হবে,
কারণ $y > x$ হলে স্বয়ংক্রিয়ভাবে~$y \nmid x$। অতএব, $x$ মৌলিক হলে
এবং কেবল তখনই সত্য সম্বন্ধ $\fn{Prime}(x)$-কে এভাবে সংজ্ঞায়িত করা যায়:
\[
\fn{Prime}(x) \defiff x \geq 2 \land \bforall{y \leq x}{(y \mid x \lif y
  = 1 \lor y = x)}
\]
এবং তাই সম্বন্ধটি আদিম পুনরাবৃত্ত।

মৌলিক সংখ্যাগুলি হলো $2$, $3$, $5$, $7$, $11$ ইত্যাদি। এই অনুক্রমের
$x$-তম মৌলিক সংখ্যাটি ফেরত দেয় এমন অপেক্ষক $p(x)$ বিবেচনা করুন;
অর্থাৎ $p(0) = 2$, $p(1) = 3$, $p(2) = 5$ ইত্যাদি। (সুবিধার জন্য
আমরা প্রায়ই $p(x)$-কে $p_x$ লিখব; অর্থাৎ $p_0=2$, $p_1=3$ ইত্যাদি।)

আমাদের কাছে যদি~$x$-এর চেয়ে বড় প্রথম মৌলিক সংখ্যাটি ফেরত দেয় এমন
অপেক্ষক $\fn{nextPrime(x)}$ থাকে, তবে আদিম পুনরাবৃত্তির সাহায্যে $p$-কে
সহজেই সংজ্ঞায়িত করা যায়:
\begin{align*}
  p(0) & = 2\\
  p(x+1) & = \fn{nextPrime}(p(x))
\end{align*}
যেহেতু $\fn{nextPrime}(x)$ হলো এমন ক্ষুদ্রতম $y$ যার জন্য $y > x$ এবং
$y$ মৌলিক, তাই সীমাহীন অনুসন্ধানের সাহায্যে একে সহজেই গণনা করা যায়।
কিন্তু ইউক্লিডের একটি ফলের কারণে সীমাবদ্ধ ন্যূনীকরণ দিয়েও একে সংজ্ঞায়িত
করা যায়: $x$ ও $\fact{x}+1$-এর মধ্যে সবসময় একটি মৌলিক সংখ্যা থাকে।
\[
  \fn{nextPrime(x)} =
  \bmin{y \leq \fact{x}+1}{(y > x \land \fn{Prime}(y))}.
\]
এ থেকে দেখা যায়, $\fn{nextPrime}(x)$ এবং তাই $p(x)$ শুধু গণনসাধ্যই
নয়, আদিম পুনরাবৃত্তও।

(কৌতূহল থাকলে ইউক্লিডের উপপাদ্যটির একটি সংক্ষিপ্ত প্রমাণ এখানে দেওয়া
হলো। ধরা যাক, $p_n$ হলো $\le x$ বৃহত্তম মৌলিক সংখ্যা এবং
$p = p_0\cdot p_1 \cdot \dots \cdot p_n$ হলো সব মৌলিক সংখ্যা~$\le x$-এর
গুণফল। হয় $p+1$ মৌলিক, নয়তো $x$ ও~$p+1$-এর মধ্যে একটি মৌলিক সংখ্যা
আছে। কেন? ধরা যাক, $p+1$ মৌলিক নয়। তাহলে কোনো মৌলিক সংখ্যা
$q \mid p+1$, যেখানে $q < p+1$। $\le x$ কোনো মৌলিক সংখ্যাই $p+1$-কে
নিঃশেষে ভাগ করে না। ($p$-এর সংজ্ঞা অনুসারে প্রত্যেক মৌলিক সংখ্যা
$p_i \le x$, $p$-কে নিঃশেষে ভাগ করে; অর্থাৎ ভাগশেষ~$0$। তাই প্রত্যেক
মৌলিক সংখ্যা $p_i \le x$ দিয়ে $p+1$ ভাগ করলে ভাগশেষ~$1$, এবং তাই
$p_i \nmid p+1$।) অতএব $q$ একটি মৌলিক সংখ্যা, যা $> x$ এবং
$< p+1$। আবার $p \le \fact{x}$, তাই $> x$ ও $\le \fact{x}+1$ এমন
একটি মৌলিক সংখ্যা আছে।)

\begin{prob}
সীমাবদ্ধ ন্যূনীকরণ ব্যবহার করে পূর্ণসংখ্যা ভাগ $d(x, y)$ সংজ্ঞায়িত করুন।
\end{prob}

\end{document}
